\documentclass[11pt,a4paper]{article}

% ── Codificación y fuentes ────────────────────────────────────────────────────
\usepackage[utf8]{inputenc}
\usepackage[T1]{fontenc}
\usepackage{lmodern}
\usepackage[spanish]{babel}

% ── Geometría y espaciado ─────────────────────────────────────────────────────
\usepackage[top=2.5cm, bottom=2.5cm, left=2.8cm, right=2.8cm]{geometry}
\usepackage{setspace}
\setstretch{1.15}
\usepackage{parskip}

% ── Tablas ────────────────────────────────────────────────────────────────────
\usepackage{booktabs}
\usepackage{tabularx}
\usepackage{array}
\usepackage{longtable}
\usepackage{enumitem}

% ── Colores y cajas ───────────────────────────────────────────────────────────
\usepackage[dvipsnames]{xcolor}
\usepackage{tcolorbox}
\tcbuselibrary{skins, breakable}

% ── Cabeceras y pies de página ────────────────────────────────────────────────
\usepackage{fancyhdr}
\usepackage{lastpage}
\pagestyle{fancy}
\fancyhf{}
\fancyhead[L]{\small\textbf{Informe de Evaluación} --- Física para Bioanalistas}
\fancyhead[R]{\small franyerica rodriguez 33278284}
\fancyfoot[C]{\small Página \thepage\ de \pageref{LastPage}}
\renewcommand{\headrulewidth}{0.4pt}

% ── Hiperenlaces ──────────────────────────────────────────────────────────────
\usepackage[colorlinks=true, linkcolor=NavyBlue, urlcolor=NavyBlue]{hyperref}

% ── Paleta de colores personalizados ─────────────────────────────────────────
\definecolor{desempeno_excelente}{HTML}{1A6B2F}
\definecolor{desempeno_bueno}{HTML}{2E5A9C}
\definecolor{desempeno_suficiente}{HTML}{8B6914}
\definecolor{desempeno_deficiente}{HTML}{C0392B}
\definecolor{desempeno_insuficiente}{HTML}{7B241C}
\definecolor{grisFondo}{HTML}{F5F5F5}
\definecolor{azulTitulo}{HTML}{1B2A6B}
\definecolor{verdeFortaleza}{HTML}{1A6B2F}
\definecolor{naranjaMejora}{HTML}{A04000}

% ── Estilos de cajas ─────────────────────────────────────────────────────────
\tcbset{
  cajaTitulo/.style={
    enhanced, breakable,
    colback=azulTitulo!8, colframe=azulTitulo,
    fonttitle=\bfseries\large, coltitle=white,
    attach boxed title to top left={yshift=-2mm, xshift=4mm},
    boxed title style={colback=azulTitulo},
    top=6pt, bottom=6pt, left=8pt, right=8pt,
  },
  cajaFortaleza/.style={
    enhanced, breakable,
    colback=verdeFortaleza!6, colframe=verdeFortaleza!60,
    fonttitle=\bfseries, coltitle=verdeFortaleza,
    top=4pt, bottom=4pt, left=8pt, right=8pt,
  },
  cajaMejora/.style={
    enhanced, breakable,
    colback=naranjaMejora!6, colframe=naranjaMejora!60,
    fonttitle=\bfseries, coltitle=naranjaMejora,
    top=4pt, bottom=4pt, left=8pt, right=8pt,
  },
  cajaRecomendacion/.style={
    enhanced, breakable,
    colback=grisFondo, colframe=gray!50,
    fonttitle=\bfseries, coltitle=black,
    top=4pt, bottom=4pt, left=8pt, right=8pt,
  },
}

% ─────────────────────────────────────────────────────────────────────────────
\begin{document}

% ══ PORTADA ══════════════════════════════════════════════════════════════════
\begin{center}
  {\Large\bfseries\color{azulTitulo} Universidad de Oriente/Núcleo Sucre --- Física para Bioanalistas}\\[4pt]
  {\large Informe de Evaluación Preliminar}\\[12pt]
  \rule{\linewidth}{1.2pt}\\[6pt]
  {\LARGE\bfseries franyerica rodriguez 33278284}\\[4pt]
  \rule{\linewidth}{0.6pt}
\end{center}

% ══ FICHA TÉCNICA ════════════════════════════════════════════════════════════
\vspace{4pt}
\begin{tcolorbox}[cajaTitulo, title=Datos del informe]
\begin{tabular}{@{}llll@{}}
  \textbf{Estudiante:}  & franyerica rodriguez 33278284  &
  \textbf{Fecha:}       & 2026-06-25 \\[4pt]
  \textbf{Archivo PDF:} & \texttt{franyerica\_rodriguez\_33278284.pdf}  &
  \textbf{Páginas:}     & 5 \\
\end{tabular}
\end{tcolorbox}

% ══ RESULTADO GLOBAL ═════════════════════════════════════════════════════════
\vspace{6pt}
\begin{center}
  \begin{tcolorbox}[
    enhanced,
    colback=white, colframe=desempeno_bueno,
    width=0.62\linewidth,
    boxrule=2pt,
    halign=center, valign=center,
    top=8pt, bottom=8pt,
  ]
    \centering
    {\large\bfseries Puntaje sugerido}\\[6pt]
    {\Huge\bfseries\color{desempeno_bueno} 8.9/10}\\[6pt]
    {\large Nivel de desempeño: \textbf{\color{desempeno_bueno} Bueno}}
  \end{tcolorbox}
\end{center}

% ══ RESUMEN GENERAL ══════════════════════════════════════════════════════════
\section*{Resumen general}
El estudiante demuestra un buen dominio de los principios de la física de fluidos aplicados al bioanálisis. La mayoría de los problemas fueron resueltos correctamente, mostrando una sólida comprensión conceptual y la habilidad para aplicar las fórmulas adecuadas. Se identificaron algunos errores menores de transcripción en una fórmula y un error de cálculo en un exponente, lo que afectó la precisión de la respuesta final en un problema.

% ══ FORTALEZAS Y ÁREAS DE MEJORA ════════════════════════════════════════════
\vspace{4pt}
\begin{minipage}[t]{0.48\linewidth}
  \begin{tcolorbox}[cajaFortaleza, title={Fortalezas}]
\begin{itemize}[leftmargin=1.5em, itemsep=2pt]
  \item Sólida comprensión de los principios de hidrostática, flotación y hemodinámica.
  \item Habilidad para aplicar correctamente las fórmulas físicas en la mayoría de los casos.
  \item Correcta identificación y manejo de datos en la mayoría de los problemas.
  \item Buen uso de unidades y consistencia en los cálculos (excepto un error en Q5).
  \item Capacidad para realizar análisis proporcionales (Ley de Poiseuille).
\end{itemize}
  \end{tcolorbox}
\end{minipage}
\hfill
\begin{minipage}[t]{0.48\linewidth}
  \begin{tcolorbox}[cajaMejora, title={Areas de mejora}]
\begin{itemize}[leftmargin=1.5em, itemsep=2pt]
  \item Precisión en la escritura de fórmulas intermedias para evitar ambigüedades o errores de transcripción.
  \item Revisión cuidadosa de los cálculos, especialmente con exponentes y potencias de diez, para asegurar la exactitud numérica.
\end{itemize}
  \end{tcolorbox}
\end{minipage}

% ══ OBSERVACIONES POR PREGUNTA ═══════════════════════════════════════════════
\section*{Observaciones por pregunta}

\begin{longtable}{>{\bfseries}p{2.8cm} p{1.6cm} p{7.0cm} p{2.8cm}}
  \toprule
  \textbf{Pregunta} & \textbf{Puntaje} & \textbf{Evaluación} & \textbf{Errores detectados} \\
  \midrule
  \endhead
  \bottomrule
  \endfoot
    \textbf{1. Presión Hidrostática y Venosa} & 2.0/2.0 & La respuesta es completamente correcta. El estudiante identificó la fórmula adecuada, sustituyó los valores correctamente y obtuvo el resultado exacto con las unidades apropiadas. & \textit{---} \\
    \midrule
    \textbf{2. Principio de Arquímedes} & 2.0/2.0 & La resolución es impecable. El estudiante calculó correctamente la fuerza de empuje, el volumen del objeto y finalmente su densidad, siguiendo todos los pasos lógicos y matemáticos de forma precisa. & \textit{---} \\
    \midrule
    \textbf{3. Hemodinámica (Continuidad y Bernoulli)} & 1.9/2.0 & El cálculo de la velocidad V2 usando la ecuación de continuidad es correcto. La aplicación de la ecuación de Bernoulli para encontrar P2 también es correcta en su ejecución numérica. Sin embargo, en la línea donde se sustituyen los valores en la ecuación de Bernoulli, se escribió incorrectamente la densidad (10600 en lugar de 1060) y el signo del término de velocidad (suma en lugar de resta). A pesar de este error de transcripción, el cálculo subsiguiente utilizó los valores y el signo correctos, lo que sugiere un descuido al escribir más que un error conceptual. & \textit{Error de transcripción en la línea de sustitución de Bernoulli: densidad incorrecta (10600 en lugar de 1060) y signo incorrecto en el término de velocidades (suma en lugar de resta), aunque el cálculo posterior fue correcto.} \\
    \midrule
    \textbf{4. Ley de Poiseuille (Análisis Proporcional)} & 2.0/2.0 & La respuesta es completamente correcta. El estudiante comprendió la relación de la Ley de Poiseuille con el radio (Q $\propto$ r\textasciicircum{}4) y calculó correctamente el porcentaje de caudal que se mantiene circulando. & \textit{---} \\
    \midrule
    \textbf{5. Flujo Viscoso y Resistencia} & 1.0/2.0 & El estudiante seleccionó la fórmula correcta de la Ley de Poiseuille para la caída de presión y extrajo los datos adecuadamente, incluyendo la conversión del diámetro a radio. Sin embargo, se cometió un error en el cálculo del denominador ($\pi$r\textasciicircum{}4), específicamente en el exponente de la potencia de diez, lo que llevó a un resultado final incorrecto en la magnitud de la presión. & \textit{Error procedimental en el cálculo del denominador ($\pi$r\textasciicircum{}4), específicamente en el exponente de la potencia de diez, lo que afectó el resultado final.} \\
    \midrule
\end{longtable}

% ══ ERRORES DETECTADOS ═══════════════════════════════════════════════════════
\section*{Análisis de errores}

\subsection*{Errores conceptuales}
\textit{Ninguno identificado.}

\subsection*{Errores procedimentales}
\begin{itemize}[leftmargin=1.5em, itemsep=2pt]
  \item Error de transcripción en una línea de la ecuación de Bernoulli (Pregunta 3).
  \item Error de cálculo en el exponente del denominador ($\pi$r\textasciicircum{}4) en la Pregunta 5.
\end{itemize}

% ══ RECOMENDACIONES AL ESTUDIANTE ════════════════════════════════════════════
\section*{Retroalimentación para el estudiante}
\begin{tcolorbox}[cajaRecomendacion, title={Dirigido a: franyerica rodriguez 33278284}]
Estimado/a estudiante, tu desempeño en este examen es bueno, demostrando una sólida base en los principios de la física de fluidos. Has aplicado correctamente la mayoría de los conceptos y fórmulas. Para mejorar aún más, te sugiero prestar especial atención a la precisión al escribir las fórmulas intermedias y al realizar cálculos con potencias de diez, ya que pequeños errores en estos aspectos pueden llevar a resultados finales incorrectos. Revisa siempre tus pasos, especialmente los exponentes, para asegurar la exactitud numérica. ¡Sigue practicando!
\end{tcolorbox}

% ══ NOTA PARA EL PROFESOR ════════════════════════════════════════════════════
\section*{Nota para el profesor}
\begin{tcolorbox}[
  enhanced, breakable,
  colback=yellow!8, colframe=yellow!60!black,
  fonttitle=\bfseries, coltitle=black,
  title={Observaciones para revisión manual},
  top=4pt, bottom=4pt, left=8pt, right=8pt,
]
El estudiante muestra un buen nivel de comprensión general. El error en la Pregunta 3 parece ser más un descuido al escribir una línea intermedia que un error conceptual profundo, ya que el cálculo final fue correcto. El error en la Pregunta 5 es de cálculo numérico (exponente), no de concepto. El puntaje sugerido refleja un buen desempeño con penalizaciones por estos errores de precisión.
\end{tcolorbox}

% ══ PIE DE INFORME ════════════════════════════════════════════════════════════
\vspace{12pt}
\begin{center}
  \footnotesize\color{gray}
  Informe generado automáticamente por el sistema de evaluación con IA (gemini-2.5-flash) y soporte LaTex. \\
  Este documento es una evaluación \textbf{preliminar} para ser revisado por el profesor antes de ser entregado al 
  estudiante. \\
  Generado el 2026-06-25 \textbullet\ BIOANALISIS Sistema Académico de Evaluación.
  \end{center}

\end{document}
